% Resume: Arumalla Mohit Krishna
% Compile with pdfLaTeX.

\documentclass[letterpaper,10pt]{article}

\usepackage[margin=0.5in]{geometry}
\usepackage[T1]{fontenc}
\usepackage{lmodern}
\usepackage{microtype}
\usepackage{titlesec}
\usepackage{enumitem}
\usepackage[hidelinks]{hyperref}
\usepackage{tabularx}
\input{glyphtounicode}

\pdfgentounicode=1
\pagestyle{empty}
\urlstyle{same}
\raggedright
\setlength{\parindent}{0pt}
\setlength{\tabcolsep}{0pt}

\titleformat{\section}
  {\large\scshape}{}{0em}{}[\titlerule]

\titlespacing*{\section}{0pt}{6pt}{3pt}

\newcommand{\entry}[2]{%
  \noindent
  \begin{tabularx}{\linewidth}
    {@{}>{\raggedright\arraybackslash}X@{\hspace{6pt}}r@{}}
    \textbf{#1} & #2
  \end{tabularx}\par
}

\newcommand{\bullets}{
  \begin{itemize}[
    leftmargin=1.2em,
    topsep=2pt,
    itemsep=1pt,
    parsep=0pt,
    partopsep=0pt
  ]
}

\newcommand{\finishbullets}{
  \end{itemize}
}

\begin{document}

%----------HEADING----------
\begin{center}
    {\LARGE\bfseries Arumalla Mohit Krishna}\\[3pt]
    +1 (631) 352-8643 $|$
    \href{mailto:mohitkrishna.arumalla@stonybrook.edu}
    {mohitkrishna.arumalla@stonybrook.edu} $|$
    \href{https://linkedin.com/in/mohit-krishna-arumalla}
    {LinkedIn}
\end{center}
\vspace{-8pt}

%-----------EDUCATION-----------
\section{Education}

\entry
  {Stony Brook University}
  {Aug 2025 -- Dec 2026}

M.S. in Computer Science \hfill New York, USA\\
\textit{
Coursework: Natural Language Processing,
Introduction to Computer Vision, Vision Language Modeling
}
\vspace{3pt}

\entry
  {Birla Institute of Technology and Science, Pilani}
  {Aug 2019 -- 2023}

B.E. in Electronics \& Communication; GPA: 8.02/10 \hfill India
%-----------PROFESSIONAL SUMMARY-----------
\section{Professional Summary}

Software engineer with 2+ years of experience building **scalable distributed services, ML infrastructure, and real-time data pipelines**. Currently pursuing an M.S. in Computer Science at Stony Brook University. Hands-on expertise in **machine learning lifecycle management, GPU model serving, and production-grade software tooling**. Skilled in **Python, Java, cloud infrastructure (AWS, GCP), Docker, and Kubernetes**, with a strong focus on robust and high-performance system design.

%-----------TECHNICAL SKILLS-----------
\section{Technical Skills}

\textbf{Languages:}
Python, Go, C++, Java, SQL\\

\textbf{ML \& Retrieval:}
PyTorch, TensorFlow/Keras, scikit-learn, NumPy, pandas,
Transformers, mBART, BGE, Vision Transformers,
Self-Supervised Learning, LoRA, Visual Prompt Tuning,
Text Embeddings, Vector Search\\

\textbf{Data \& Systems:}
gRPC, Kafka, Apache Flink, Flink SQL, Debezium/CDC,
PostgreSQL, Redis, Elasticsearch, Airflow, BigQuery\\

\textbf{Infrastructure:}
NVIDIA Triton, AWS, GCP, Docker, Kubernetes, Linux, Git,
CI/CD, Istio/Envoy, New Relic

%-----------EXPERIENCE-----------
\section{Experience}

\entry
  {\textbf{Software Engineer II} $|$ \textbf{Infra.Market}, Bangalore}
  {Apr 2024 -- Aug 2025}

\bullets

  \item Built **semantic product search infrastructure** integrating Amazon Titan text embeddings and Elasticsearch vector search for intuitive natural-language product discovery through chat interfaces.

  \item Engineered a **distributed CDC pipeline** using Flink, Kafka, Debezium, PostgreSQL, RocksDB, and Elasticsearch; reduced synchronization latency from 10–15 minutes to 500 ms with exactly-once processing and fault tolerance.

  \item Deployed Istio/Envoy service mesh across 15+ microservices with **distributed tracing (New Relic)**, enhancing production monitoring, debugging, and reliability in a dynamic environment.

\finishbullets

\entry
  {\textbf{Software Engineer I} $|$ \textbf{Infra.Market}, Chennai}
  {Jul 2023 -- Mar 2024}

\bullets

  \item Owned Search and Location microservices; built high-throughput **gRPC services handling 500K+ daily requests** for product discovery and geolocation, ensuring reliable system operations under load.

  \item Identified and fixed PostgreSQL connection pool leaks using Gorm ORM, reducing connection exhaustion by 85\%; validated improvements with **concurrent load testing at 10K RPS** to ensure robustness.

  \item Developed **Airflow data pipelines** for BigQuery-to-Elasticsearch migration with optimized query execution, dynamic filtering, and automated monitoring to maintain up-to-date search indices.

\finishbullets

%-----------PROJECTS-----------
\section{Projects}

\entry
  {GPU-Optimized Personalized Product Search}
  {Aug 2026 -- Oct 2026}

\textit{Python, PyTorch, BGE, NVIDIA Triton, Go, gRPC}

\bullets

  \item Developed two-tower retrieval model with
  \texttt{BAAI/bge-small-en-v1.5} query/product embeddings
  (384 dimensions) and a trainable user tower combining query
  embeddings with interaction history.

  \item Built an inference pipeline featuring precomputed product embeddings,
  approximate nearest-neighbor search, and a Go/gRPC gateway;
  integrated NVIDIA Triton for efficient GPU serving of query encoders.

  \item Designed batching and concurrency benchmarks measuring latency,
  throughput, GPU utilization, and cost-efficiency; evaluated
  recall and ranking metrics under variable traffic loads.

\finishbullets

\entry
  {Attention-Based Hallucination Detection in NMT}
  {Apr 2026 -- May 2026}

\textit{Python, PyTorch, mBART, Transformers, scikit-learn}

\bullets

  \item Built an English-German hallucination-detection pipeline leveraging
  teacher-forced mBART attention features from WMT22/WMT23 quality
  estimation datasets.

  \item Extracted comprehensive cross-/self-attention metrics across decoder layers and heads,
  including entropy, attention mass statistics, variance, and grounding ratios.

  \item Trained token-, sentence-, and head-level classifiers achieving 0.712 F1 and 0.695 ROC-AUC;
  proposed hallucination-aware reranking combining localization, risk, and attention signals.

\finishbullets

\entry
  {Prompt/LoRA for Self-Supervised Learning}
  {Ongoing}

\textit{Python, PyTorch, Vision Transformers}

\bullets

  \item Investigating domain-specific visual prompting and LoRA adaptation
  in a self-supervised ViT framework to learn transferable
  representations with parameter efficiency across image domains.

  \item Evaluating performance trade-offs of backbone models without prompts/LoRA;
  designing experiments to measure classification accuracy and robustness under domain shift.

\finishbullets

\end{document}